%% kapa: exam=1 round=31 year=2020 period=1 form=A subject=real_estate subject_ko=부동산학원론 seq=4 number=84
%% source: Q-Net 감정평가사 기출문제 (cst00309) artlSeq=5207219 file=extracted/1008/5207219/A형/A형/2020년도 제31회 감정평가사 1차 시험_1교시_A형.pdf
%% conversion: kordoc 4.24.51 -> tools/gen1.py (auto); review=proofread-sonnet; flags=none
\begin{question}{84}
( )에 들어갈 내용으로 옳은 것은?
\begin{framedtext}
○ ( )(이)란 임장활동의 전 단계 활동으로 여러 가지 물적 증거를 수집하고 탁상 위에서 검토하여 1차적으로 하자의 유무를 발견하려는 작업이다. \par
○ ( )의 과정은 위험사례를 미리 발견하기 위한 노력 또는 그 기초 작업이다. \par
\end{framedtext}
\begin{choices}
  \item 보정
  \item 심사
  \item 판독
  \item 면책사항
  \item 권리보증
\end{choices}
\end{question}
\answer{3}
